\section{Warmup: Approximation algorithms for the coverage problem} \label{sec:coverage}

In this section, we consider the \coverage problem with heterogeneous costs. According to the preprocessing, we assume that interface costs are in $\sbr{0,1}$ and $\OPT \geq 1/2$. 
The idea behind our ILP formulation and the algorithms is inspired by the rounding procedure known for \textsc{Set Cover} \cite{vazirani2001approximation} and related problems \cite{Angelidakis2018ShortestPQ}.


\subsection{ILP formulations of the \coverage problem}

We start with the ILP formulation of the \coverage problem. Let $x_{i, v}$ be a binary variable that is $1$ iff the interface $i\in\lambda\br{v}$ is assigned to vertex $v\in V$. If a given vertex $v$ is a cheap, then without large cost increase we can enforce $x_{i, v}=1$, for every $i\in\lambda\br{v}$. Otherwise, if a vertex $v$ is expensive, we can assume that in an optimal solution we have $\sum_{i\in \lambda\br{v}} x_{i, v}\cdot c\br{i, v} \geq 1$, which corresponds to the fact that $\COST\br{\lambda_A, v}\geq 1$. Thus, we can write the following ILP formulation of the \coverage problem:
\begin{align}
    \text{min} \quad & M \nonumber\\
    \text{s.t.} \quad & \sum_{i\in\lambda\br{v}}c\br{i, v}\cdot x_{i, v} \leq M \quad \forall v\in V \label{eqt:covering_cost}\\
    & \sum_{i\in\lambda\br{u}\cap\lambda\br{v}}\min\brc{x_{i, u}, x_{i, v}} \geq 1 \quad \forall uv\in E \label{eqt:coverage}\\
    &\begin{cases}
         x_{i, v} = 1 \quad \forall i\in\lambda\br{v} & \text{if } \sum_{i\in \lambda\br{v}} c\br{i, v} < 1\\
        \sum_{i\in \lambda\br{v}} c\br{i, v}\cdot x_{i, v} \geq 1 & \text{otherwise}
    \end{cases}\quad \forall v\in V \label{eqt:costs_coverage}\\
    & x_{i, v} \in \{0, 1\} \quad \forall v\in V, i\in\lambda\br{v} \label{eqt:integrality_coverage}
\end{align}
Constraints~\eqref{eqt:covering_cost} enforce that the cost of each vertex does not exceed the global cost $M$. Constraints~\eqref{eqt:coverage} ensure that for every edge $uv\in E\br{G}$, there exists a common activated interface. Constraints~\eqref{eqt:costs_coverage} encode the assumption that for every cheap vertex, all of its interfaces are activated, while for every expensive vertex, the cost of activated interfaces is at least $1$. Finally, constraints~\eqref{eqt:integrality_coverage} require the solution to be integral.
Note that the above formulation is not linear due to constraints \eqref{eqt:coverage}. However, by adding additional variables $z_{i, uv}$, we can linearize those by replacing them with the constraints $\sum_{i\in\lambda\br{u}\cap\lambda\br{v}}z_{i, uv} \geq 1$ together with $z_{i, uv} \leq x_{i, u}$ and $z_{i, uv} \leq x_{i, v}$ for every $i\in\lambda\br{u}\cap\lambda\br{v}$ and $uv\in E$. Thus, we can solve the LP relaxation of this formulation in polynomial time.

% This ILP formulation can easily be adapted to \textsc{min-sum} \coverage by removing \eqref{eqt:covering_cost} and replacing the objective function by $\sum_{v \in V} \sum_{i\in\lambda\br{v}} c\br{i, v} \cdot x_{i, v}$.


\subsection{A $k$-approximation algorithm for the Coverage problem}
In this section we show that a simple rounding of the LP solution can be used to obtain a $k$-approximation. Note that for this simple algorithm the preprocessing is not necessary, which makes it easily adaptable for the min-sum variant of the \coverage problem (assuming the ILP is appropriately changed).

\begin{algorithm}[H]
\caption{$k$-approximation algorithm} \label{alg:2}%A relaxation-based $k$-approximation algorithm.
Solve the LP relaxation and obtain a fractional solution $\{x_{i, v}\}_{i\in\lambda(v), v\in V}$.

Set $\lambda_A(v) \gets \{i \in \lambda(v) \colon x_{i, v} \geq \frac{1}{k}\}$, for every $v \in V$.

\Return{$\{\lambda_A(v)\}_{v\in V}$.}
\end{algorithm}

\begin{theorem}
\Cref{alg:2} is a $k$-approximation algorithm for the \coverage problem.
\end{theorem}

\begin{proof}
Consider an edge $uv\in E$. By averaging, there exists an interface $i\in\lambda\br{u}\cap\lambda\br{v}$ such that $\min\brc{x_{i, u}, x_{i, v}} \geq \frac{1}{k}$. Thus, $i\in\lambda_A\br{u}\cap\lambda_A\br{v}$ so $\lambda_A$ is a covering assignment. Moreover, since $x_{i, v} \geq \frac{1}{k}$ for every $i\in\lambda_A\br{v}$, we have that $c\br{\lambda_A\br{v}} = \sum_{i\in\lambda_A\br{v}}c\br{i, v} \leq k \cdot \sum_{i\in\lambda_A\br{v}}c\br{i, v} \cdot x_{i, v}$. Thus, we have $\COST\br{\lambda_A} \leq k \cdot \max_{v\in V}\sum_{i\in\lambda\br{v}} c\br{i, v} \cdot x_{i, v} \leq k \cdot \OPT$.
\end{proof}


\subsection{An $O\br{\log n}$-approximation algorithm for the \coverage problem}
In order to extract an integral solution with quality independent of $k$, we employ a randomized procedure. For each interface type $i\in\sbr{k}$, we draw a uniformly random threshold $t_i \in \sbr{0, 1}$ and activate the interface at a vertex if the corresponding variable scaled up by a factor of $O\br{\log n}$ exceeds $t_i$. The procedure is formalized by the following pseudocode.

\begin{algorithm}[H]
\caption{$O(\log n)$-approximation algorithm for the \coverage problem} \label{alg:1}%A relaxation-based $O(\log n)$-approximation algorithm.

Solve the LP relaxation and obtain a fractional solution $\{x_{i, v}\}_{i\in\lambda(v), v\in V}$.

Pick independent uniformly random thresholds $t_i \in \sbr{0, 1}$, for each $i \in \sbr{k}$.

$\lambda_A(v) \gets \{i \in \lambda(v) \colon 2\ln m\cdot x_{i, v} \geq t_i\}$, for every $v \in V$.

\Return{$\{\lambda_A(v)\}_{v\in V}$.}
\end{algorithm}

\begin{theorem}
\Cref{alg:1} is an $O(\log n)$-approximation algorithm for the \coverage problem which succeeds with high probability.
\end{theorem}
\begin{proof}
Firstly, we claim that the algorithm indeed returns a feasible solution:
    \begin{lemma}
With high probability, \Cref{alg:1} returns an assignment $\lambda_A$, such that for every edge $uv \in E$, $\lambda_A(u) \cap \lambda_A(v) \neq \emptyset$.
    \end{lemma}
\begin{proof}
    Consider an edge $uv\in E$. We have:
\begin{align*}
    \Pr\sbr{\lambda_A\br{u}\cap\lambda_A\br{v}=\emptyset} & = \prod_{i\in\lambda\br{u}\cap\lambda\br{v}}\Pr\sbr{t_i > 2\ln m \cdot \min\brc{x_{i, u}, x_{i, v}}}\\
    & = \prod_{i\in\lambda\br{u}\cap\lambda\br{v}}\br{1-2\ln m \cdot\min\brc{x_{i, u}, x_{i, v}}}\\&
    \leq \prod_{i\in\lambda\br{u}\cap\lambda\br{v}}\e{-2\ln m \cdot\min\brc{x_{i, u}, x_{i, v}}} 
    \\& = \e{-2\ln m \cdot\sum_{i\in\lambda\br{u}\cap\lambda\br{v}}\min\brc{x_{i, u}, x_{i, v}}}\\
    & \leq \e{-2\ln m} = \frac{1}{m^2}
\end{align*}

By applying the union bound over all $m$ edges we get that the algorithm returns a covering assignment with probability at least $1-\frac{m}{m^2} = 1-\frac{1}{m}$.
\end{proof}

Now, we also wish to show that the approximation factor achieved by our algorithm is upper bounded by $O\br{\log n}$. We have the following lemma:
\begin{lemma}
    \Cref{alg:1} returns a solution of cost at most $12\ln n \cdot \OPT$ with high probability.
\end{lemma}
\begin{proof}
Fix an expenssive vertex $v\in V$. Otherwise, if $v$ is cheap, then the cost of activating all interfaces at $v$ is at most $1\leq2\cdot\OPT$. We want to show that with high probability, the cost of $v$ is at most $12\ln n$ times the cost of $v$ in the fractional solution. 
By $X_{v, i}$ denote the random variable that is $1$ if $2\ln m\cdot x_{i, v} \geq t_i$ and $0$ otherwise. Then we can write $\COST\br{v, \lambda_A} = \sum_{i\in\lambda\br{v}}c\br{i, v} \cdot X_{v, i}$. We have: $\mathbb{E}\sbr{X_{v, i}} = \Pr\sbr{i \in \lambda_A\br{v}} = 2\ln m\cdot x_{i, v}$. Thus, $\mathbb{E}\sbr{\COST\br{v, \lambda_A}} = 2\ln m \cdot \sum_{i\in\lambda\br{v}}c\br{i, v} \cdot x_{i, v}$. Applying the Chernoff bound with $\delta = 2$ we get 
\begin{align*}
    \Pr\sbr{\COST\br{v, \lambda_A} \geq 3\cdot \mathbb{E}\sbr{\COST\br{v, \lambda_A}}} & 
    \leq \br{\frac{e^2}{27}}^{\mathbb{E}\sbr{\COST\br{v, \lambda_A}}}
    \leq \e{-\mathbb{E}\sbr{\COST\br{v, \lambda_A}}}\leq \frac{1}{m^2}
\end{align*}
where the last inequality is by using the fact that, $\sum_{i\in\lambda\br{v}}c\br{i, v} \cdot x_{i, v} \geq 1$. By applying the union bound over all $n$ vertices we get that with probability at least $1-\frac{n}{m^2}\geq 1-\frac{m+1}{m^2}$ for every expensive $v\in V$, $\COST\br{v, \lambda_A} \leq 6\ln m \cdot \sum_{i\in\lambda\br{v}}c\br{i, v} \cdot x_{i, v}$. Thus, with high probability we have $\COST\br{\lambda_A} \leq 6\ln m \cdot \max_{v\in V}\sum_{i\in\lambda\br{v}}c\br{i, v} \cdot x_{i, v} \leq 6\ln m \cdot \OPT\leq 12\ln n \cdot \OPT$. 
\end{proof}
Combining the above lemmas gives the desired result.
\end{proof}

\section{Approximation algorithm for the \connectivity problem} \label{sec:connectivity}

In this section we move towards the \connectivity problem with heterogeneous costs. We again assume that the instance is preprocessed, so we have that all costs are in $\sbr{0,1}$ and $\OPT \geq 1/2$.

\subsection{ILP formulation of the \connectivity problem}

We use a similar ILP to model the \connectivity problem, however we use additional flow variables to encode the connectivity requirement. The $x_{i, v}$ variables are defined as before. For each edge $e$, we add a variable $y_{e}$ which is $1$ iff the edge $e$ is covered by the solution. For any $S\subseteq V$, let $\delta(S)$ denote the set of edges with exactly one endpoint in $S$. Global connectivity is enforced via cut constraints: for every non-trivial cut $(S, V\setminus S)$, there needs to be at least one covered edge crossing the cut. Let $y\br{\delta\br{S}}=\sum_{e \in \delta(S)} y_e$. We can formulate the \connectivity problem as follows:
\begin{align}
    \text{min} \quad & M \nonumber\\
    \text{s.t.} \quad & \sum_{i\in\lambda\br{v}}c\br{i, v} \cdot x_{i, v} \leq M \quad \forall v\in V \label{eqt:4}\\
    & \sum_{i\in\lambda\br{u}\cap\lambda\br{v}}\min\brc{x_{i, u}, x_{i, v}} \geq y_{uv} \quad \forall uv\in E\\
    &\begin{cases}
            x_{i, v} = 1 \quad \forall i\in\lambda\br{v} & \text{if } \sum_{i\in \lambda\br{v}} c\br{i, v} < 1\\
            \sum_{i\in \lambda\br{v}} c\br{i, v}\cdot x_{i, v} \geq 1 & \text{otherwise}
    \end{cases} \quad \forall v\in V \label{eqt:5}\\
    & y\br{\delta\br{S}}\geq 1 \quad \forall S\subset V\colon S\neq\emptyset,\, S\neq V \label{eqt:6}\\
    & x_{i, v}  \in \{0, 1\} \quad \forall v\in V, i\in\lambda\br{v} \nonumber\\
    & y_{uv} \in \{0, 1\} \quad \forall uv\in E \label{eqt:integrality_connectivity}
\end{align}

Observe that the above ILP has an exponential number of constraints, but we can solve its LP relaxation in polynomial time using the ellipsoid method and a separation oracle for the cut constraints, which can be implemented using a max-flow algorithm.

\subsection{An $O(\log^2 n)$-approximation algorithm for the \connectivity problem}
Obtaining an integral solution for the \connectivity problem requires a more careful rounding procedure than for the \coverage problem. We wish to ensure that at least one connected spanning subgraph is covered by the assignment returned by the algorithm. To do so, we start by randomly sampling edges according to the $y$-variables of the fractional solution, in a way so that with high probability, the sampled subgraph $H$ is connected. To achieve this goal, we set the probability of choosing $e$ to be $p_e=\min\{1, 5\ln m \cdot y_e\}$. Then, the goal of the algorithm becomes covering $H$. To do so, we repeat the following process: we draw uniformly random thresholds $t_i$ for each type of interface $i\in \sbr{k}$, and activate $i$ at any vertex for which $5\ln m \cdot x_{i,v}\geq t_i$. Intuitively, each such iteration costs $O\br{\log n\cdot \OPT}$ in expectation and covers $\Omega\br{1}$ fraction of edges of $H$. After repeating this for a fixed number of $T = O\br{\log n}$ rounds, it can be shown that the algorithm covers all edges of $H$ almost surely. The formal description of the algorithm is given by the following pseudocode.

\begin{algorithm}[H]
\caption{$O\br{\log^2n}$-approximation algorithm for the \connectivity problem} \label{alg:3} %A relaxation-based $O(\log^2 n)$-approximation algorithm.

Solve the LP relaxation and obtain a fractional solution $\{x_{i, v}\}_{i\in\lambda(v), v\in V}, \{y_{uv}\}_{uv\in E}$.

For every edge $e\in E$, $p_e \gets \min\{1, 5\ln m \cdot y_e\}$. 

With probability $p_e$ set $Y_e \gets 1$, otherwise set $Y_e \gets 0$. 

Let $H$ be the subgraph of $G$ induced by the edges $e$ such that $Y_e = 1$.

$T \gets \frac{2\ln m}{1-1/e}$.

$\lambda_A\gets \{\emptyset\}_{v\in V}$.

\For{$j \in \sbr{T}$}{
Pick independent uniformly random thresholds $t_i \in (0, 1)$, for each $i \in \sbr{k}$.

$\lambda_A^j(v) \gets \{i \in \lambda(v) \colon 5\ln m \cdot x_{i, v} \geq t_i\}$, for every $v \in V$.

$\lambda_A \gets \lambda_A^j \cup \lambda_A$.
}
\Return{$\lambda_A$}
\end{algorithm}

\begin{theorem}
    \Cref{alg:3} is an $O\br{\log^2 n}$-approximation algorithm for the \connectivity problem which succeeds with high probability.
\end{theorem}
\begin{proof}
    In order to show that the algorithm returns a connecting assignment, we first show that with high probability, the subgraph $H$ is connected. It should be remarked that constructing $H$ explicitly is not necessary for the algorithm to work. Nonetheless, we do so for the sake of convenience of our analysis. We start with the following lemma:
    \begin{lemma}
        With high probability, $H$ is connected.
    \end{lemma}
    \begin{proof}
        Fix any $S\subset V$ with $S\neq\emptyset$ and $S\neq V$. By definition of the LP, there exists a natural number $\ell\in\mathbb{N}$, such that $\ell\leq y\br{\delta_H\br{S}}<\ell+1$. Let $\delta_H\br{S}=E\br{H}\cap \delta\br{S}$.
        We have:
        \begin{align*}
            \Pr\sbr{\spr{\delta_H\br{S}}=0} & = \prod_{e\in\delta\br{S}}\Pr\sbr{Y_e=0} = \prod_{e\in\delta\br{S}}(1-p_e) \\&\leq \prod_{e\in\delta\br{S}}\e{-p_e} \leq \e{-\sum_{e\in\delta(S)}\min\brc{1,\, 5\ln m\cdot y_e}}
            \\& \leq \e{-5\ell\ln m} = \frac{1}{m^{5\ell}}
        \end{align*}
        % where the last inequality uses the LP constraint $y\br{\delta\br{S}}\geq 1$ for every non-trivial cut, so $\sum_{e\in\delta(S)} y_e \geq \ell$.

        A cut in $G$ is called an $\alpha$-minimum cut if its value is less than $\alpha$ times the value of a minimum cut in $G$. 
        We now make use of the following Karger's lemma \cite{MinimumCutsInNearLinearTime, RandomSamplingInCutFlowAndNetworkDesignProblems}:
        \begin{lemma}[Karger's lemma]
            The number of $\alpha$-minimum cuts in a graph is at most $n^{2\alpha}$.
        \end{lemma}
        Observe that by construction of the LP, every non-trivial fractional cut has value at least $1$, so there are at most $n^{2\ell+2}$ such cuts that have fractional value less than $\ell+1$. Applying the union bound over all of them we obtain:
\begin{align*}
    \Pr\sbr{\exists\, S\subset V,\; S\neq\emptyset,\, S\neq V\colon \spr{\delta_H\br{S}}=0}
    &\leq \sum_{\ell=1}^{\infty}\sum_{\substack{S\subset V,\,S\neq\emptyset,\,S\neq V,\\ y\br{\delta_H\br{S}}<\ell+1}}\Pr\sbr{\spr{\delta_H\br{S}}=0} \\
    &\leq \sum_{\ell=1}^{\infty}n^{2\ell+2}\cdot \frac{1}{m^{5\ell}} \leq \sum_{\ell=1}^{\infty}\frac{1}{m^{\ell}} = \frac{1}{m-1}
\end{align*}
where the last inequality is due to the fact, that assuming that $G$ is not a tree, we get that $m\geq n$. Otherwise, \coverage and \connectivity are equivalent and one can use the $O\br{\log n}$-approximation algorithm instead.

Thus, with probability at least $1-\frac{1}{m-1}$, every non-trivial cut $S$ of $G$, has at least one edge of $H$ crossing it, which means that $H$ spans all vertices of $G$.
\end{proof}

Having shown that the subgraph $H$ is connected, we now wish to show that the assignment $\lambda_A$ returned by the algorithm covers all of the edges of $H$ almost surely. 
Let $j$ be an index of some iteration of the for loop. Let $H_j$ be the (possibly empty) subgraph of $H$ that is not covered by $\lambda_A$ at the beginning of the $j$-th iteration of the for loop. We wish to show, that in expectation, the solution $\lambda_A$ produced in the $j$-th iteration of the for loop, $\lambda_A^j$, covers a significant fraction of the edges in $H_j$. We have the following lemma:
\begin{lemma}
    We have
    $\mathbb{E}\sbr{\spr{\cov\br{\lambda_A^j, H_j}}} \geq \br{1-1/e} \cdot \mathbb{E}\sbr{\spr{E\br{H_j}}}$.
\end{lemma}
\begin{proof}
    Consider any edge $uv\in E\br{G}$. Observe that $\mathbb{E}\sbr{Y_{uv}}\leq 5\ln m \cdot y_{uv}$. By definition of the LP we get that:
    \begin{align*}
        \sum_{i\in\lambda\br{u}\cap\lambda\br{v}}\min\brc{x_{i, u}, x_{i, v}} & \geq y_{uv} \geq \frac{\mathbb{E}\sbr{Y_{uv}}}{5\ln m}\cdot
    \end{align*}
    Therefore, the probability that $uv$ is uncovered by $\lambda_A^j$ is at most:
    \begin{align*}
        \Pr\sbr{\lambda_A\br{u}^j\cap\lambda_A\br{v}^j=\emptyset} & = \prod_{i\in\lambda\br{u}\cap\lambda\br{v}}\Pr\sbr{t_i > 5\ln m\cdot\min\brc{x_{i, u}, x_{i, v}}} \\ 
        & = \prod_{i\in\lambda\br{u}\cap\lambda\br{v}}\br{1-5\ln m\cdot\min\brc{x_{i, u}, x_{i, v}}} \\& \leq \prod_{i\in\lambda\br{u}\cap\lambda\br{v}}\e{-5\ln m\cdot\min\brc{x_{i, u}, x_{i, v}}}\\ 
        & = \e{-5\ln m\cdot\sum_{i\in\lambda\br{u}\cap\lambda\br{v}}\min\brc{x_{i, u}, x_{i, v}}} \\& \leq \e{-5\ln m\cdot\frac{\mathbb{E}\sbr{Y_{uv}}}{5\ln m}}=\e{-\mathbb{E}\sbr{Y_{uv}}}
    \end{align*}
    
    Thus, we have $\Pr\sbr{\lambda_A\br{u}^j\cap\lambda_A\br{v}^j\neq\emptyset}\geq 1-\e{-\mathbb{E}\sbr{Y_{uv}} }$. By summing up over all edges in $H_j$ we lower bound the expected number of edges covered by $\lambda_A^j$ as follows:
    \begin{align*}
        \mathbb{E}\sbr{\spr{\cov\br{\lambda_A^j, H_j}}} & = \sum_{uv\in E\br{H_j}}\Pr\sbr{\lambda_A\br{u}^j\cap\lambda_A\br{v}^j\neq\emptyset}  
        \geq \sum_{uv\in E\br{H_j}}\br{1-\e{-\mathbb{E}\sbr{Y_{uv}} }}
        \\& \geq \br{1-1/e}\cdot \sum_{uv\in E\br{H_j}}{\mathbb{E}\sbr{Y_{uv}}} \geq \br{1-1/e} \cdot \mathbb{E}\sbr{\spr{E\br{H_j}}}
    \end{align*}
    where the second inequality is due to the fact that $1-e^{-x}\geq (1-1/e)\cdot x$, for $x\in\sbr{0,1}$ and the last inequality is obtained by linearity of expectation.
\end{proof}

Let $I$ be the random variable denoting the number of iterations needed to cover all edges of $H$, i.e., the smallest $j$ such that $\cov\br{\lambda_A, H} = E\br{H}$ after $j$ iterations (assume that the number of iterations might eventually get larger then $T$). We show that with high probability $I \leq T$, so the algorithm returns a feasible solution.
\begin{lemma}
    \label{lem:expected-iterations}
    $\Pr\sbr{I > T} = O\br{\frac{1}{m}}$.
\end{lemma}
\begin{proof}
    Let $\alpha = 1-1/e$. By definition, $\spr{E\br{H_{j+1}}}=\spr{E\br{H_{j}}}-\spr{\cov\br{\lambda_A^j, H_j}}$. By linearity of expectation and the previous lemma, $\mathbb{E}\sbr{\spr{E\br{H_{j+1}}}} \leq \br{1-\alpha}\cdot \mathbb{E}\sbr{\spr{E\br{H_{j}}}}$. Unfolding the recurrence gives $\mathbb{E}\sbr{\spr{E\br{H_{j}}}} \leq m e^{-\alpha(j-1)}$. By Markov's inequality, $\Pr\sbr{\spr{E\br{H_{j}}}\geq 1} \leq m e^{-\alpha(j-1)}$. Since $T = \frac{2\ln m}{1-1/e} = \frac{2\ln m}{\alpha}$, setting $j = T$ gives:
    $$\Pr\sbr{I > T} = \Pr\sbr{\spr{E\br{H_{T}}}\geq 1} \leq m e^{-\alpha(T-1)} \leq m e^{-2\ln m+\alpha} = O\br{\frac{1}{m}}.$$
    since $e^{\alpha}=O\br{1}$.
\end{proof}

It remains to show that the cost of the solution returned by the algorithm is at most $O\br{\log^2 n}\cdot \OPT$. We have with the following lemma:
\begin{lemma}
    With high probability, for every vertex $v\in V\br{G}$ and every iteration $j \in \sbr{T}$, $\COST\br{v, \lambda_A^j} \leq 15\ln m \cdot \OPT$.
\end{lemma}
\begin{proof}
Fix a vertex $v\in V$ and an iteration $j\in\sbr{T}$. If $v$ is cheap, then its overall cost is at most $1\leq 2\cdot\OPT$ regardless of activated interfaces. Thus, assume $v$ to be expensive, so $\sum_{i\in\lambda\br{v}}c\br{i,v} \cdot x_{i,v} \geq 1$.
By $X_{v,i}^j$ denote the indicator that is true iff $i\in \lambda_A^j$. Then, $\COST\br{v, \lambda_A^j} = \sum_{i\in\lambda\br{v}}c\br{i,v} \cdot X_{v,i}^j $. We have $\mathbb{E}\sbr{X_{v,i}^j} = 5\ln m\cdot x_{i,v}$, hence $\mathbb{E}\sbr{\COST\br{v,\lambda_A^j}} = 5\ln m\cdot \sum_{i\in\lambda\br{v}}c\br{i,v} \cdot x_{i,v}$. Applying the Chernoff bound with $\delta=2$ gives:
\begin{align*}
    \Pr\sbr{\COST\br{v, \lambda_A^j} \geq 3\cdot\mathbb{E}\sbr{\COST\br{v,\lambda_A^j}}}
    \leq \br{\frac{e^2}{27}}^{\mathbb{E}\sbr{\COST\br{v,\lambda_A^j}}}
    \leq e^{-\mathbb{E}\sbr{\COST\br{v,\lambda_A^j}}}
    \leq \frac{1}{m^5},
\end{align*}
where the last inequality uses $\mathbb{E}\sbr{\COST\br{v,\lambda_A^j}} = 5\ln m\cdot\sum_{i\in\lambda\br{v}}c\br{i,v} \cdot x_{i,v} \geq 5\ln m$. Hence $\COST\br{v,\lambda_A^j} \leq 15\ln m\cdot\sum_{i\in\lambda\br{v}} c\br{i,v} \cdot x_{i,v} \leq 15\ln m\cdot\OPT$ with probability at least $1-\frac{1}{m^5}$.

Applying the union bound over all $n \leq m + 1$ vertices and all $T \leq \frac{2\ln m}{1-1/e}$ iterations, the above holds for all $v\in V$ and $j\in\sbr{T}$ with probability at least $1 - \frac{T\cdot \br{m+1}}{m^5} = 1 - O\br{\frac{1}{m^3}}$.
\end{proof}
Consequently, summing up over all $T$ iterations, with high probability:
$$\COST\br{\lambda_A} \leq \sum_{j=1}^{T}\max_{v\in V\br{G}}\brc{\COST\br{v, \lambda_A^j}} \leq T\cdot 15\ln m\cdot\OPT = \frac{30\ln^2 m}{1-1/e}\cdot\OPT\leq \frac{120\ln^2 n}{1-1/e}\cdot\OPT$$

Combining the above lemmas gives the desired result.
\end{proof}

